\newpage

\section{Metodología de ensayo y validación}
La metodología de ensayo propuesta tiene como objetivo verificar el cumplimiento de las especificaciones funcionales y de desempeño del generador de funciones/arbitrario, en el dominio temporal. Se definen el instrumental, las condiciones de medición, los ensayos y los criterios de aceptación.

\subsection{Instrumental y condiciones de medición}
Los ensayos se realizarán bajo condiciones ambientales de laboratorio. Se utilizará el siguiente instrumental:
\begin{itemize}
	\item Osciloscopio con ancho de banda adecuado (idealmente $\geq 5\times$ la frecuencia máxima a evaluar), sondas pasivas y entrada con terminación de 50\,$\Omega$ o alta impedancia según se disponga.
	\item Analizador de espectro (si se dispone) o FFT del osciloscopio para evaluación espectral.
	\item Multímetro TRMS para verificación de amplitud en baja frecuencia.
	\item Cables coaxiales BNC de baja longitud y adaptadores/terminadores adecuados.
\end{itemize}

\subsection{Ensayos funcionales}
Se verificará el correcto funcionamiento de las funciones principales:
\begin{itemize}
	\item \textbf{Generación de formas de onda:} El usuario deberá seleccionar la función de \textit{Tipo de señal} que permitirá cambiar entre las formas de ondas almacenadas previamente senoidal, cuadrada, triangular/rampa, pulso, electrocardiograma, diente de sierra, gausseana, etc.
	\item \textbf{Configuración de parámetros:} El usuario debe seleccionar el parámetro a configurar según corresponda frecuencia, amplitud, offset DC, duty cycle y fase. Con el encoder le permitirá ir variando el valor de cada uno y mediante botones de selección podrá cambiar el multiplicador.
	\item \textbf{Visualización de Configuraciones:} definida a través de la interfaz de visualización que se use (display o PC).
	\item Verificación de interfaz de control (USB/serial/Ethernet) y respuesta a comandos.
\end{itemize}

El resultado satisfactorio de esta etapa radica en la posibilidad de que el usuario/desarrollador pueda manejar de forma correcta las configuraciones mencionadas y que estas se vean claramente representadas en la interfaz de visualización elegida.

%\subsection{Ensayos de frecuencia y resolución}
%Se evaluará la exactitud y resolución de frecuencia mediante puntos de prueba representativos (por ejemplo 1\,kHz, 10\,kHz, 100\,kHz, 300 KHz, 500 KHz, 700 KHz, 1\,MHz, 5\,MHz, 10\,MHz y frecuencia máxima del equipo):
%
%\begin{itemize}
%%	\item \textbf{Error de frecuencia:} $\Delta f = f_\text{medida} - f_\text{programada}$ (si es posible en base a la calidad del instrumental usado).
%	\item \textbf{Resolución efectiva:} menor cambio programable observable en la salida (igual que item anterior).
%\end{itemize}
%
%Como criterio de aceptación se adoptará un error máximo permitido en frecuencia (en ppm o en \%), definido en función de los objetivos del proyecto y del oscilador de referencia utilizado.

\subsection{Ensayos de tipo de señal, frecuencia, amplitud y offset}
En esta sección de presentan todos los tipos de ensayos que van a realizarse, con el fin de darle la mayor caracterización al equipo que se va a realizar. Cada ensayo debe hacerse para cada una de las señales convencionales (senoidal, cuadrada, triangular) y además para alguna arbitraria cargada con una cantidad de muestras definida.

\subsubsection{Ensayo de frecuencia}
Para realizar este ensayo el usuario deberá seleccionar primeramente el tipo de canal que quiere configurar y presionar OK. Luego mediante el encoder rotativo y los botones de cambio de multiplicador, deberá moverlos de acuerdo a la frecuencia que quiera a salida según la tabla proporcionada.

\begin{table}[H]
	\centering
	\label{tab:freq_set_meas}
	\begin{tabular}{|l|c|r|l|}
		\hline
		Tipo de se\~nal & $f_{set}$ & $f_{med}$ & Observaciones \\ \hline
		& 1\,Hz      &  & \\ \hline
		& 10\,Hz     &  & \\ \hline
		& 100\,Hz    &  & \\ \hline
		& 1\,kHz     &  & \\ \hline
		& 10\,kHz    &  & \\ \hline
		& 100\,kHz   &  & \\ \hline
		& 1\,MHz     &  & \\ \hline
		& 3\,MHz     &  & \\ \hline
		& 5\,MHz     &  & \\ \hline
		& 10\,MHz    &  & \\ \hline
		& 15\,MHz    &  & \\ \hline
		& 20\,MHz    &  & \\ \hline
		& 25\,MHz    &  & \\ \hline
	\end{tabular}
	\caption{\textbf{Ensayo frecuencia:} frecuencia configurada vs frecuencia medida ($V_{pp}$ de 20\,V, $V_{off}=0$).}
\end{table}

\subsubsection{Ensayo de amplitud}
Para dicho ensayo el usuario primeramente deberá seleccionar el canal que quiere cambiar la amplitud de salida en Volts de pico a pico o Volts de pico. Una vez seleccionada deberá hacer uso del encoder y botones de cambio de multiplicador para poder ajustar el valor final definido por la siguiente tabla.

\begin{table}[H]
	\centering
	\label{tab:amp_set_meas}
	\begin{tabular}{|l|r|r|r|l|}
		\hline
		Tipo de se\~nal & $V_{pp,set}$ & $V_{pp,med}$ & $V_{off,med}$ & Observaciones \\ \hline
		& 0.2\,V  &  &  & \\ \hline
		& 0.5\,V  &  &  & \\ \hline
		& 1.0\,V  &  &  & \\ \hline
		& 2.0\,V  &  &  & \\ \hline
		& 5.0\,V  &  &  & \\ \hline
		& 10.0\,V &  &  & \\ \hline
		& 15.0\,V &  &  & \\ \hline
		& 20.0\,V &  &  & \\ \hline
	\end{tabular}
	\caption{\textbf{Ensayo amplitud:} amplitud configurada vs amplitud medida (frecuencia 10\,kHz, $V_{off}=0$).}
\end{table}

\subsubsection{Ensayo de offset}
El procedimiento es el mismo que los ensayos anteriores, seleccionando en este caso el boton de cambio de offset o señal de DC.

\begin{table}[H]
	\centering
	\label{tab:offset_set_meas}
	\begin{tabular}{|l|r|r|r|l|}
		\hline
		Tipo de se\~nal & $V_{off,set}$ & $V_{off,med}$ & $V_{pp,med}$ & Observaciones \\ \hline
		& -5.0\,V &  &  & \\ \hline
		& -4.0\,V &  &  & \\ \hline
		& -3.0\,V &  &  & \\ \hline
		& -2.0\,V &  &  & \\ \hline
		& -1.0\,V &  &  & \\ \hline
		& -0.5\,V &  &  & \\ \hline
		& 0.0\,V  &  &  & \\ \hline
		& +0.5\,V &  &  & \\ \hline
		& +1.0\,V &  &  & \\ \hline
		& +2.0\,V &  &  & \\ \hline
		& +3.0\,V &  &  & \\ \hline
		& +4.0\,V &  &  & \\ \hline
		& +5.0\,V &  &  & \\ \hline
	\end{tabular}
	\caption{\textbf{Ensayo offset:} offset configurado vs offset medido (frecuencia 10\,kHz, $V_{pp}$ fijo).}
\end{table}

\subsubsection{Ensayo de planitud}
A diferencia de los ensayos anteriores el usuario ahora deberá configurar tanto la amplitud en un valor definido (preferentemente el máximo) e ir incrementando la frecuencia y visualizar la tensión de salida máxima. Esto se hace con el fin de poder ver que tan estable es la salida frente a diferentes frecuencias y a su vez las limitaciones propias del equipo. Además al hacerlo tanto para alta impedancia como baja impedancia las condiciones serán distintas ya que en baja impedancia la exigencia será mucho mayor.

\begin{table}[H]
	\centering
	\label{tab:sine_flatness_basic}
	\begin{tabular}{|c |r| l|}
		\hline
		$f$ & $V_{pp,med}$ & Condición (Hi-Z / 50$\Omega$) \\
		\hline
		1\,kHz    &  & \\ \hline
		10\,kHz   &  & \\ \hline
		100\,kHz  &  & \\ \hline
		1\,MHz    &  & \\ \hline
		5\,MHz    &  & \\ \hline
		10\,MHz   &  & \\ \hline
		20\,MHz   &  & \\
		\hline
		25\,MHz   &  & \\
		\hline
	\end{tabular}
	\caption{\textbf{Ensayo senoidal:} amplitud medida vs frecuencia.}
\end{table}

\subsubsection{Ensayos específicos}
Incluyen todos los ensayos que depende exclusivamente del tipo de señal.

\begin{table}[H]
	\centering
	\label{tab:square_duty_rise}
	\begin{tabular}{|c| r| r| r| c| r| r| l|}
		\hline
		$f_{set}$ & $f_{med}$ &
		$V_{pp,set}$ & $V_{pp,med}$ &
		Duty$_{set}$ & Duty$_{med}$ &
		$t_r$ / $t_f$ &
		Obs. \\
		\hline
		1\,kHz   &  & 1.0\,V &  & 50\% &  &  & \\ \hline
		100\,kHz &  & 1.0\,V &  & 50\% &  &  & \\ \hline
		1\,MHz   &  & 1.0\,V &  & 50\% &  &  & \\
		\hline
	\end{tabular}
	\caption{Ensayo específico: señal cuadrada (duty y tiempos).}
\end{table}


%\subsection{Integridad de señal en el dominio temporal}
%Para cada forma de onda, especialmente a frecuencias elevadas, se evaluará:
%
%\begin{itemize}
%	\item \textbf{Para senoidal:} ruido observable, estabilidad y distorsión.
%	\item \textbf{Para cuadrada/pulso:} tiempos de subida y bajada (10--90\%), overshoot/undershoot para ciertas frecuencias seleccionadas, duty real versus duty programado.
%\end{itemize}
%
%Dado que se trata de un proyecto final a nivel educativo y no un dispositivo comercial, la validación de la etapa será en función de los resultados obtenidos.

%\subsection{Ensayos espectrales (armónicos, THD, SFDR)}
%Se caracterizará la calidad espectral mediante FFT del osciloscopio o analizador de espectro. Para una senoidal (por ejemplo 1\,MHz y 10\,MHz) se evaluará:
%
%\begin{itemize}
%	\item Niveles de armónicos (2º, 3º, etc.).
%	\item THD o THD+N (si el instrumento lo permite).
%	\item SFDR: diferencia entre el tono fundamental y el mayor espurio.
%\end{itemize}
%
%Se verificará que el contenido espectral sea consistente con la forma cargada y que no existan espurios dominantes no esperados. En caso de aparecer mencionarlos y dar una breve justificación.